\FloatBarrier
\section{Methodology}
%------------------------------------------------------------------------------
\subsection{Dataset and Preprocessing}

The IMDb Large Movie Review Dataset~\cite{maas2011} (50,000
reviews, balanced binary labels) was split 80/20 train-test
with \texttt{random\_state=42}, and 10\% of
training data was withheld as a validation set. Reviews were
cleaned by removing HTML tags and non-alphabetic characters,
then lowercased. A vocabulary of the 5,000 most frequent
training-set tokens was used; all sequences were truncated or
zero-padded to 200 tokens (Table~\ref{tab:hyperparams}).

\subsection{Network Architecture}

All four experiments share the same model structure. The
only thing that changes within each experiment is the choice of
Act$_1$ and Act$_2$ in the two dense layers. The shared architecture
is: an embedding layer
(vocabulary 5,000, dimension 128), a bidirectional recurrent
Layer~1 (64 units per direction), a recurrent Layer~2 (32
units per direction), Dropout(0.4), Dense Layer~1 (64 units)
with Act$_1$, Dense Layer~2 (32 units) with Act$_2$, and a
single output unit trained with \texttt{BCEWithLogitsLoss}.
The four experiments differ only in the recurrent cell type
(LSTM, GRU) and whether Layer 2
is bidirectional (Table~\ref{tab:architectures}).

\begin{table}[htbp]
  \caption{Experimental Architecture Summary}
  \label{tab:architectures}
  \centering
  \footnotesize
  \setlength{\tabcolsep}{4pt}
  \renewcommand{\arraystretch}{1.08}
  \begin{tabular}{cllr}
    \toprule
    Exp. & Layer 1 & Layer 2 & Params \\
    \midrule
    1 & BiLSTM(64) & LSTM(32)    & ${\approx}$764K \\
    2 & BiGRU(64)  & GRU(32)     & ${\approx}$734K \\
    3 & BiLSTM(64) & BiLSTM(32)  & ${\approx}$787K \\
    4 & BiGRU(64)  & BiGRU(32)   & ${\approx}$752K \\
    \bottomrule
  \end{tabular}
\end{table}

\subsection{Activation Functions Under Investigation}

Ten activation functions were evaluated
(Table~\ref{tab:activations}). The classical family is
ReLU~\cite{nair2010}, LeakyReLU~\cite{maas2013},
PReLU~\cite{he2015}, ELU~\cite{clevert2015}, Tanh and
SELU~\cite{klambauer2017}. The modern smooth family is
GELU~\cite{hendrycks2016}, SiLU (Swish)~\cite{ramachandran2017},
Mish~\cite{misra2020} and Hardswish~\cite{howard2019}. The latter
four had not previously been benchmarked in gated recurrent dense
layers for sentiment classification.

\begin{table}[htbp]
  \caption{Activation Functions: Families and Key Properties}
  \label{tab:activations}
  \centering\footnotesize
  \setlength{\tabcolsep}{3pt}
  \renewcommand{\arraystretch}{1.05}
  \resizebox{\columnwidth}{!}{%
  \begin{tabular}{lll}
    \toprule
    Function & Type & Formula / Key Property \\
    \midrule
    ReLU      & Classical & $\max(0,x)$; dead neurons for $x{<}0$ \\
    ELU       & Classical & $\alpha(e^x{-}1)$ for $x{\leq}0$; mean-centering \\
    LeakyReLU & Classical & $0.01x$ for $x{<}0$; low dead-neuron risk \\
    Tanh      & Classical & Bounded; saturates for $|x|\gg1$ \\
    GELU      & Modern & $x\cdot\Phi(x)$; probabilistic gating \\
    SiLU      & Modern & $x\cdot\sigma(x)$; self-gating \\
    Mish      & Modern & $x\cdot\tanh(\text{softplus}(x))$; smooth \\
    SELU      & Modern & Self-normalising; implicit batch norm \\
    PReLU     & Modern & Learnable slope $a$; input-dependent \\
    Hardswish & Modern & Piecewise; dead zone for $x{<}{-3}$ \\
    \bottomrule
  \end{tabular}%
  }
\end{table}

\subsection{Training Configuration and Evaluation}

Every setting in Table~\ref{tab:hyperparams} was held the same
across the 2{,}000 runs of the primary sweep, so the activation
pair is the only thing that changes between models. The
architecture parameters (vocabulary, embedding, layer sizes,
dropout) replicate Ahmed et al.~\cite{ahmed2025} exactly. The
training parameters (optimizer, learning rate, batch size,
scheduler, clipping) were chosen by us because Ahmed et al. did
not report them. After training, accuracy, precision, recall, and
F1-score were measured on the held-out 10,000-review test set.

The early-learning follow-up of
Section~\ref{sec:convergence_protocol} is a separate experiment
and departs from this protocol in three documented respects.

\subsection{Replication Protocol}
\label{sec:replication}

In the primary sweep, each of the $10\times10\times4 = 400$
configurations was trained under five seeds,
$\{10,20,30,40,50\}$ (Table~\ref{tab:hyperparams}), giving
$2{,}000$ models. The seed is applied immediately before the
model and its data loaders are constructed, setting
\texttt{random}, \texttt{numpy}, \texttt{torch} and
\texttt{torch.cuda}, with \texttt{cudnn.deterministic=True} and
\texttt{cudnn.benchmark=False} to pin kernel selection as well.
A given (pair, seed) is therefore exactly reproducible.

The split is deliberately \emph{not} a function of the seed: it
stays pinned at \texttt{random\_state=42}. The variance we report
is therefore the variance of retraining the same model on the same
data. That is the quantity a single-run comparison assumes to be
negligible. Split variance would only inflate it, so these
estimates are conservative.

Runs were executed five at a time, one worker per seed. Each
worker is a separate process with its own RNG and deterministic
kernel selection, so concurrency cannot perturb a result; we
confirmed this by re-running sampled configurations serially and
recovering identical metrics.

\begin{table}[htbp]
  \caption{Shared Hyperparameter Configuration}
  \label{tab:hyperparams}
  \centering\footnotesize
  \setlength{\tabcolsep}{3pt}
  \renewcommand{\arraystretch}{1.05}
  \resizebox{\columnwidth}{!}{%
  \begin{tabular}{lll}
    \toprule
    Hyperparameter & Value & Source \\
    \midrule
    Vocab / Seq.\ length / Embed. & 5{,}000 / 200 / 128 & \cite{ahmed2025} \\
    RNN Layer 1 / Layer 2 units   & 64 / 32 per direction & \cite{ahmed2025} \\
    Dense Layer 1 / Layer 2 units & 64 / 32 & \cite{ahmed2025} \\
    Dropout rate                  & 0.4 & \cite{ahmed2025} \\
    Batch size / Max epochs       & 64 / 10 & Ours \\
    Early stopping patience       & 3 epochs (val.\ loss) & Ours \\
    Optimiser                     & Adam, $\text{lr}=10^{-3}$ & Ours \\
    LR scheduler                  & ReduceLROnPlateau ($\times$0.5) & Ours \\
    Loss / Gradient clipping      & BCEWithLogitsLoss / $L_2\leq1.0$ & Ours \\
    Weight initialisation         & Kaiming uniform (Linear) & PyTorch default \\
    Data split seed               & 42, fixed for all runs & Ours \\
    Training seeds (sweep)        & 10, 20, 30, 40, 50 & Ours \\
    Training seeds (follow-up)    & Sec.~\ref{sec:convergence_protocol} & Ours \\
    \bottomrule
  \end{tabular}%
  }
\end{table}

\subsection{Early-Learning Follow-Up Protocol}
\label{sec:convergence_protocol}

The early-learning experiment reported in
Section~\ref{sec:learning} is a separate, smaller study. It does
not inherit the primary sweep's protocol, and we state the
differences here because two of them affect how its numbers
should be read.

\textit{Scope.} It trains BiGRU+GRU only, and only the ten
configurations that use the same function in both dense layers
(X+X), rather than all 100 ordered pairs. The dataset, split,
vocabulary, architecture sizes, optimiser, scheduler, batch size,
maximum epochs and early-stopping rule are all as in
Table~\ref{tab:hyperparams}.

\textit{Measurement.} It records validation loss and validation
accuracy after every epoch. The test set is never touched. Its
accuracies are therefore validation accuracies on the held-out
10\% of the training split, and are not comparable with the test
accuracies reported elsewhere in this paper.

\textit{Seeding.} This experiment sets \texttt{torch.manual\_seed}
and \texttt{numpy.random.seed} only, and leaves the cuDNN flags at
their library defaults. The narrower scope is deliberate: it keeps
these runs bit-compatible with the earlier single-seed convergence
result, and enabling \texttt{cudnn.deterministic} changes the
numerics enough to flip epoch-one outcomes. That sensitivity is
itself informative. Epoch one sits close to a threshold, in which
a model either leaves the chance-level plateau within one pass or
does not, so small numerical differences can decide the epoch-one
figure while leaving later epochs unaffected. These runs are
reproducible on a fixed machine and library version, but we do not
claim reproducibility across different cuDNN kernel selections.

\textit{Seed sets.} The follow-up was run twice, under two
disjoint five-seed sets, $\{42,0,1,2,3\}$ and
$\{10,20,30,40,50\}$, for 100 runs in total. The first set
contains seed 42, which is also the seed of the single-run
ranking it is compared against, so on its own it cannot serve as
an independent check of that ranking. The second set shares no
seed with the first and matches the seeds of the primary sweep,
which allows the epoch-one ranking to be tested on runs that have
no seed in common with those used to establish it.

\subsection{On Reproducing the Baseline of Ahmed et al.}
\label{sec:baseline_repro}

Ahmed et al.~\cite{ahmed2025} describe their architecture in full
and we copy it exactly, but they do not report how they trained
it: optimizer, learning rate, batch size, epochs, early stopping,
seed and data split are all missing. Exact reproduction is
therefore impossible, and we adopt the fully documented settings of
Table~\ref{tab:hyperparams} instead.

Replication sharpens what can be said about their 88.40\%. Our best
configuration averages 87.81\% (Table~\ref{tab:top_configs})
with a 95\% upper bound of 88.01\%,
and no run among the 2{,}000 reached 88.40\% (maximum 88.39\%), so
the benchmark sits outside the interval of our strongest setting.
Since they report a single run without a seed or standard
deviation, and we measure the per-cell standard deviation of this
model family at 0.38 points, their figure is best read as one draw
of unknown uncertainty. We treat it as a reference point, not a
target.
